%=====================================================================
% 01 — TÓM TẮT VI + EN — WP6 compressed
%=====================================================================

\noindent\textbf{\textit{Tóm tắt:}}

\noindent
\textbf{Quản lý doanh nghiệp khoa học và công nghệ (DN KH\&CN) và doanh nghiệp khởi nghiệp sáng tạo (DN KNST) chủ yếu thuộc thẩm quyền cấp tỉnh nhưng cần liên thông và tổng hợp ở cấp trung ương. Theo nghiên cứu khoa học thiết kế (\textit{Design Science Research}, DSR), nghiên cứu đề xuất một kiến trúc tham chiếu phần mềm (SRA) tỉnh--trung ương có khả năng truy vết. Phát hiện trung tâm là quyền ghi cần được biểu diễn như một gán có phạm vi thẩm quyền và thời gian hiệu lực, không phải một thuộc tính địa bàn tĩnh, với một máy trạng thái năm bước bảo toàn bất biến một writer hợp lệ tại một thời điểm (single-active-writer). SRA tách lõi ổn định khỏi các điểm biến thiên, xác định nguồn dữ liệu theo nhóm thuộc tính, đồng thời quản trị cấu hình quy trình và liên thông mà không khóa cấu trúc lưu trữ vật lý. Mười kịch bản trước triển khai phát hiện hai khoảng trống: chuyển quyền ghi theo thời gian và khả năng cùng tồn tại của các phiên bản khác nhau. Sau tinh chỉnh, cùng rubric cho kết quả 2 kịch bản đáp ứng trực tiếp, 8 đáp ứng có điều kiện và không còn rủi ro kiến trúc chưa được hấp thụ ở mức L2. Các kết quả được cô đọng thành ba nguyên tắc thiết kế có khả năng chuyển giao cho các nền tảng chính phủ đa cấp khác: phân rã thẩm quyền theo nhóm thuộc tính--phạm vi--thời gian; tách lõi pháp lý ổn định khỏi biến thể vận hành; và duy trì provenance/liên thông độc lập với topology lưu trữ. Đây là kiểm chứng kiến trúc dựa trên kịch bản, không phải bằng chứng thực nghiệm; nguyên mẫu vẫn cần kiểm thử chuyển quyền, nâng cấp, liên thông, phục hồi, an toàn và các thuộc tính chất lượng đo được.}

\vspace{0.4em}
\noindent\textbf{\textit{Từ khóa:}} \textbf{doanh nghiệp khoa học và công nghệ, kiến trúc dữ liệu, kiến trúc tham chiếu, khởi nghiệp sáng tạo, liên thông hệ thống, quy trình cấu hình.}

\vspace{0.2em}
\noindent\textbf{\textit{Chỉ số phân loại:}} \textbf{1.2, 2.2}

\newpage
\begin{english}
\begin{center}
{\bfseries A provincial--central reference architecture for a digital platform managing science and technology enterprises and innovative startups\par}
\vspace{0.4em}
{\bfseries Dinh Quang Nguyen\textsuperscript{1}, Huy Cuong Nguyen\textsuperscript{2*}\par}
\vspace{0.3em}
{\itshape \textsuperscript{1}Faculty of Information Technology, Posts and Telecommunications Institute of Technology, 96A Tran Phu Street, Ha Dong Ward, Hanoi, Vietnam\par}
\vspace{0.15em}
{\itshape \textsuperscript{2}Technology and Science Enterprise Development Center, National Agency for Technology Entrepreneurship and Commercialization, Ministry of Science and Technology, 115 Tran Duy Hung Street, Yen Hoa Ward, Hanoi, Vietnam\par}
\end{center}

\noindent\textbf{\textit{Abstract:}}

\noindent
\textbf{Science and technology enterprises (STEs) and innovative startups are primarily administered provincially but require interoperability with central systems and national aggregation. Following Design Science Research (DSR), this study develops a traceable provincial--central software reference architecture (SRA). The central finding is that write authority must be represented as an assignment scoped by jurisdiction and effective time, not a static territorial attribute, governed by a five-state transfer machine that preserves a single-active-writer invariant. The SRA separates a stable core from variation points, assigns authoritative sources by attribute group, and governs workflow configuration and interoperability without fixing storage topology. Ten pre-implementation scenarios exposed two gaps: temporal transfer of write authority and coexistence of different operational versions. After refinement, the same rubric yielded two directly supported scenarios, eight conditionally supported scenarios and no L2-unabsorbed architectural risks. The findings are distilled into three transferable design principles, applicable to other multi-tier government platforms: decompose authority by attribute group, scope and effective time; separate a stable legal core from controlled operational variation; and preserve provenance/interoperability independently of storage topology. These results are scenario-based architecture verification rather than empirical system validation; prototype tests are still required for authority transfer, version-skew upgrades, interoperability, recovery, security and measurable quality attributes.}

\vspace{0.4em}
\noindent\textbf{\textit{Keywords:}} \textbf{configurable workflow, data authority, innovative startups, interoperability, reference architecture, science and technology enterprises.}

\vspace{0.2em}
\noindent\textbf{\textit{Classification numbers:}} \textbf{1.2, 2.2}
\end{english}

\newpage
